\documentclass[11pt,a4paper]{article}

% ── Packages ──────────────────────────────────────────────────────────────────
\usepackage[T1]{fontenc}
\usepackage[utf8]{inputenc}
\usepackage{lmodern}
\usepackage[margin=1in]{geometry}
\usepackage{hyperref}
\usepackage{amsmath,amssymb}
\usepackage{booktabs}
\usepackage{enumitem}
\usepackage{listings}
\usepackage{xcolor}
\usepackage{microtype}
\usepackage{titlesec}
\usepackage{fancyhdr}
\usepackage{csquotes}
\usepackage{parskip}
\usepackage{mdframed}

% ── Colours ───────────────────────────────────────────────────────────────────
\definecolor{codebg}{HTML}{F5F5F5}
\definecolor{codecomment}{HTML}{6A737D}
\definecolor{codestring}{HTML}{032F62}
\definecolor{codekw}{HTML}{D73A49}
\definecolor{accent}{HTML}{0366D6}
\definecolor{boxbg}{HTML}{FFF8E1}
\definecolor{boxborder}{HTML}{F9A825}

% ── Listings ──────────────────────────────────────────────────────────────────
\lstset{
  backgroundcolor=\color{codebg},
  basicstyle=\ttfamily\small,
  commentstyle=\color{codecomment},
  keywordstyle=\color{codekw}\bfseries,
  stringstyle=\color{codestring},
  breaklines=true,
  frame=single,
  framesep=4pt,
  rulecolor=\color{gray!40},
  xleftmargin=8pt,
  xrightmargin=8pt,
  captionpos=b,
  numbers=left,
  numberstyle=\tiny\color{gray},
  numbersep=6pt,
}

% ── Callout box ───────────────────────────────────────────────────────────────
\newmdenv[
  backgroundcolor=boxbg,
  linecolor=boxborder,
  linewidth=2pt,
  leftline=true, rightline=false, topline=false, bottomline=false,
  innerleftmargin=10pt,
  innerrightmargin=10pt,
  innertopmargin=6pt,
  innerbottommargin=6pt,
  skipabove=\medskipamount,
  skipbelow=\medskipamount,
]{callout}

% ── Section formatting ────────────────────────────────────────────────────────
\titleformat{\section}{\large\bfseries}{\thesection.}{0.5em}{}[\titlerule]
\titleformat{\subsection}{\normalsize\bfseries}{\thesubsection.}{0.5em}{}

% ── Header / footer ──────────────────────────────────────────────────────────
\pagestyle{fancyplain}
\fancyhf{}
\lhead{\textit{Working Backwards in Code}}
\rhead{\thepage}
\renewcommand{\headrulewidth}{0.4pt}

% ── Hyperref ─────────────────────────────────────────────────────────────────
\hypersetup{
  colorlinks=true,
  linkcolor=accent,
  urlcolor=accent,
  citecolor=accent,
  pdftitle={Working Backwards in Code: Explore, Plan, Write Agents as the Computational Press Release},
  pdfauthor={Sean Chatman},
}

% ── Macros ────────────────────────────────────────────────────────────────────
\newcommand{\code}[1]{\texttt{#1}}
\newcommand{\pr}{\textsc{pr}}
\newcommand{\gap}{\Delta}

% ═════════════════════════════════════════════════════════════════════════════
\begin{document}
% ═════════════════════════════════════════════════════════════════════════════

\begin{titlepage}
  \centering
  \vspace*{2cm}
  {\LARGE\bfseries Working Backwards in Code\\[0.5em]
   Explore, Plan, Write Agents\\[0.3em]
   as the Computational Press Release\par}
  \vspace{2cm}
  {\large Sean Chatman\\[0.3em]
   \textit{ChatmanGPT}\par}
  \vspace{0.8cm}
  {\normalsize \today\par}
  \vspace{3cm}
  \begin{abstract}
    \noindent
    Amazon's Working Backwards process begins with the press
    release---a description of the finished product written before a line of
    code exists.  The press release is not a lie.  It is a design instrument:
    a future state held in present tense to pull a team toward it.  This paper
    argues that frontier software development is structurally identical to
    Working Backwards, and that Claude Code's Explore/Plan/Write agent
    architecture is its computational implementation.  The \emph{explore agent}
    measures the gap between press release and reality.  The \emph{plan agent}
    prioritises which part of the gap to close next.  The \emph{write agent}
    closes it.  The one-task-at-a-time sequential cycle fails frontier work
    not because tasks are hard but because it destroys the press release:
    each completed ticket becomes its own end, severed from the original
    vision.  We demonstrate with the \code{pictl}/\code{bcinr} arc, where an
    ambitious process-mining registry and a branchless algorithm library were
    each press releases written before their implementations existed---and
    where the agent architecture was the mechanism for measuring and closing
    the gap between aspiration and ground truth.
  \end{abstract}
\end{titlepage}

\tableofcontents
\newpage

% ─────────────────────────────────────────────────────────────────────────────
\section{Is the Press Release a Lie?}
% ─────────────────────────────────────────────────────────────────────────────

Amazon's Working Backwards process requires teams to write the press release
\emph{first}---before the product spec, before the architecture review, before
any code.  The press release describes the product as if it already ships.  It
uses present tense.  It quotes imaginary customers.  It states capabilities
that do not yet exist.

Is it wrong?  Is it dishonest?  Is it a lie?

No.  The press release is a \textbf{design instrument}.  Its purpose is not to
describe what exists but to make vivid what should exist, so that everyone
building toward it shares the same target.  The press release is a commitment
device operating on desired reality, not a report operating on current reality.

Now consider the \code{pictl} algorithm registry.

The registry claimed 41 algorithms: DFG, Alpha++, Inductive Miner, ILP, Genetic
Algorithm, 36 more.  It declared output types, quality scores, speed scores,
deployment profiles.  Most of these algorithms were stubs or had architectural
defects that prevented them from running on real data.

Is the registry a lie?

Only if you read it as documentation.  If you read it as a press release---a
description of the intended final system, written in code before the
implementations existed---then it is exactly correct.  It is Working Backwards
applied to a software registry: \emph{this is the system we are building
toward}.

The error was not writing the press release.  The error was losing track of
which parts were press release and which were production---and having no
systematic mechanism for measuring the gap between the two.

\begin{callout}
  \textbf{Central claim of this paper:} The purpose of the Explore/Plan/Write
  agent architecture is to hold the press release (plan), measure the gap
  between press release and reality (explore), and close it incrementally
  (write).  The one-task-at-a-time sequential cycle fails frontier work
  because it orphans each task from the press release, turning means into ends.
\end{callout}

% ─────────────────────────────────────────────────────────────────────────────
\section{Working Backwards: The Press Release as Engineering Primitive}
% ─────────────────────────────────────────────────────────────────────────────

\subsection{The Structure of Working Backwards}

Amazon's Working Backwards process has four artefacts, produced in this order:

\begin{enumerate}
  \item \textbf{Press Release (\pr):} The product as it will be described on
    launch day.  Written in plain language.  Includes the headline, the
    problem solved, the key features, and a customer quote.  Present tense.
    Future state.

  \item \textbf{FAQ:} The hard questions the press release glosses over.
    What happens when X fails?  How does it handle Y at scale?  What does
    it cost?

  \item \textbf{User Experience Narrative:} A day-in-the-life story walking
    through the product from a user's perspective, in enough detail to expose
    design gaps.

  \item \textbf{Visual / Mock:} Screenshots, diagrams, or prototypes showing
    what the product looks like---before building it.
\end{enumerate}

Only after all four are complete does implementation begin.  The press release
does not change during implementation unless the product fundamentally changes.
It is the fixed star that implementation navigates by.

\subsection{The Computational Equivalent}

In a software system, the press release is expressed not in English but in
types, APIs, registries, and configuration schemas.  The \code{pictl} registry
was a press release written in TypeScript:

\begin{lstlisting}
{
  id: 'inductive_miner',
  name: 'Inductive Miner',
  outputType: 'tree',
  quality: { fitness: 0.90, precision: 0.85 },
  speed: 30,
}
\end{lstlisting}

This says: \emph{we are building a system that includes an inductive miner with
these properties}.  It is a design commitment, not an implementation claim.
The FAQ for this press release would include: ``What does the WASM function
actually return?''  The UX narrative would include: ``A user calls
\code{discover\_inductive\_miner} and receives a process tree.''  The visual
would be the output schema.

None of these existed.  That is the gap.  The gap is not evidence that the
press release was wrong.  It is the measurement of how much work remains.

\subsection{Why Frontier Work Requires the Press Release First}

Conventional software development can work bottom-up: implement a module,
test it, add it to the system.  This works because the problem domain is
understood.  The shape of the solution is known before implementation begins.

Frontier work---new algorithms, new runtime targets, new architectural
patterns---cannot work bottom-up.  The shape of the solution is precisely what
is unknown.  Bottom-up exploration produces a collection of implemented pieces
with no coherent target.

The press release solves this by establishing the target \emph{before} the
shape of the solution is known.  Implementation proceeds top-down from the
press release, not bottom-up from components.  The question is never ``what
should I implement next?'' but ``what is the largest gap between the press
release and reality, and how do I close it?''

% ─────────────────────────────────────────────────────────────────────────────
\section{The Agent Architecture as Working Backwards Engine}
% ─────────────────────────────────────────────────────────────────────────────

Claude Code's Explore/Plan/Write agent architecture maps precisely onto the
Working Backwards process.

\begin{center}
\begin{tabular}{@{}lll@{}}
  \toprule
  \textbf{Working Backwards} & \textbf{Agent Type} & \textbf{Output} \\
  \midrule
  Press Release (desired state)     & Plan Agent    & Plan file \\
  Gap measurement (FAQ, UX)         & Explore Agent & Signal (facts about reality) \\
  Implementation (close the gap)    & Write Agent   & State change (code, tests) \\
  \bottomrule
\end{tabular}
\end{center}

\subsection{The Plan File Is the Press Release}

The plan file---a single Markdown document updated across sessions---is the
computational press release.  It is:

\begin{itemize}
  \item Written before implementation, describing the desired end state
  \item Revised when the product fundamentally changes (not when a task is hard)
  \item The fixed star that write agents navigate by
  \item Visible to the human before any implementation begins
\end{itemize}

When the \code{pictl} arc reached the conclusion that 0 of 41 algorithms were
production-ready, the plan file was not amended with a note.  It was
\textbf{rewritten}---from ``audit and fix 41 algorithms'' to ``implement 12
branchless families in \code{bcinr}''.  This is exactly what Working Backwards
prescribes: when the product changes, the press release changes.  Everything
else follows from the new press release.

\subsection{Explore Agents Measure the Gap}

An explore agent's job is to produce a precise measurement of the gap
$\gap = \text{press release} - \text{reality}$.

In the \code{pictl} arc:

\begin{itemize}
  \item Press release: 41 algorithms, all callable, all returning typed outputs,
    all measurable for fitness, precision, generalization, simplicity.
  \item Reality (discovered by explore agents): 0 production-ready, 24 with
    mutex panics, 2 returning wrong output types, 2 stubs calling DFG
    internally.
  \item Gap: $\gap = 41$ algorithms deep, architectural root cause is
    \code{RefCell<Mutex>} in single-threaded WASM.
\end{itemize}

The gap measurement did not make the press release wrong.  It made the gap
\emph{legible}.  A legible gap can be evaluated: is it worth closing, or should
the press release be revised to a closer horizon?

\subsection{Write Agents Close the Gap}

Write agents are gap-closure engines.  They receive a specific gap---``the
\code{mask} family has 3 functions that need SIMD-safe implementations''---and
produce code that closes it.  They do not decide what to implement.  They do
not evaluate whether the press release is correct.  They close the gap they
were given.

This division of labour is critical.  An agent that simultaneously evaluates
the press release and implements toward it will collapse the two roles: it will
implement what it can implement rather than what the press release requires.
This is how registries get populated with stubs.

% ─────────────────────────────────────────────────────────────────────────────
\section{Why the Sequential Task Cycle Destroys the Press Release}
% ─────────────────────────────────────────────────────────────────────────────

\subsection{The Ticket as the Atom of Sequential Cycles}

The dominant model of software development treats the \emph{ticket} as the
atomic unit of work.  A ticket has a title, a description, acceptance criteria,
and a status: To Do, In Progress, Done.  Tickets are created, assigned, worked,
and closed.  Velocity is measured in tickets per sprint.

This model is excellent for known work.  It is catastrophic for frontier work.

The ticket model has one implicit contract: \textbf{tickets are completable}.
Every ticket has a Done state that can be reached by a competent engineer in a
bounded amount of time.  The system has no representation for a ticket that
\emph{should not be completed}---a ticket that is, in Working Backwards terms,
a gap that the press release revision has made irrelevant.

\subsection{Ticket Completion Severs Tasks from the Press Release}

Consider the \code{pictl} inductive miner task.  As a ticket, it reads:

\begin{lstlisting}
Title: Implement discover_inductive_miner
Status: In Progress
Acceptance: Function is callable, returns output
\end{lstlisting}

A developer under time pressure closes this ticket by implementing a stub:
\code{discover\_inductive\_miner} calls \code{discover\_dfg} and returns its
output.  The acceptance criteria are technically met.  The ticket is marked
Done.

Is this wrong?  In the ticket model: no.  The function is callable.  It returns
output.  Done.

In the Working Backwards model: yes.  The press release says the inductive
miner returns a \emph{process tree}, not a DFG.  The stub closes the ticket
while widening the gap between press release and reality.  The Done state of
the ticket has become an obstacle to the press release.

This is the core failure mode: \textbf{ticket Done is not press release Done}.
A ticket can be marked Done while the gap to the press release grows.  Over
enough sprints, the registry fills with Done tickets and a press release that
is further from reality than when the project started.

\subsection{Sequential Cycles Measure Velocity, Not Gap}

Sprint velocity measures tickets closed per unit time.  It is an excellent
metric for known work, where each ticket is an independent unit of value.

For frontier work, velocity is worse than useless.  It measures the rate at
which the team is producing Done tickets---which, as shown above, can be
entirely decorrelated from the rate at which the gap to the press release is
closing.

A team that closes 10 stub tickets per sprint is producing high velocity and
negative progress.  The gap is growing.  The press release is receding.
But the burndown chart looks excellent.

The correct metric for frontier work is $\Delta\gap$ per unit time: the rate
of gap closure toward the press release.  This metric requires knowing both
the press release (the desired state) and the current reality (the measured
state).  The ticket model provides neither.

\subsection{Sequential Cycles Cannot Detect Architectural Gaps}

The deepest problem with sequential task cycles is that they explore the
system one task at a time.  If task $i$ fails, the team moves to task $i+1$.
The failure of task $i$ is treated as a local defect, not as a signal about
the architecture.

In the \code{pictl} arc, sequential debugging would have found:
\begin{itemize}
  \item Task 1: \code{alpha\_plus\_plus} crashes. Fix: unclear. Status: blocked.
  \item Task 2: \code{genetic\_algorithm} crashes. Fix: unclear. Status: blocked.
  \item Task 3: \code{aco} crashes. Fix: unclear. Status: blocked.
\end{itemize}

Three sprints in, three blocked tickets.  The architectural conclusion---that
\code{RefCell<Mutex>} is incompatible with single-threaded WASM---is not
visible.  It requires seeing all three failures simultaneously and asking:
what do these have in common?

Parallel exploration---launching explore agents against all 41 algorithms
simultaneously---produced this architectural conclusion in a single session.
The pattern was immediately legible because all the data arrived together.

\textbf{Sequential cycles surface local facts.  Parallel exploration surfaces
architectural facts.  Architectural facts are the only ones that can revise the
press release.}

% ─────────────────────────────────────────────────────────────────────────────
\section{The pictl Arc as Working Backwards in Practice}
% ─────────────────────────────────────────────────────────────────────────────

\subsection{The Press Release: 41 Algorithms}

The \code{pictl} registry was a press release written in TypeScript.  It
described a process mining platform with 41 algorithms spanning discovery,
conformance, ML analysis, streaming, import/export, and autonomic control.
Each algorithm had metadata: speed, quality, output type, deployment profiles.

This was not a lie.  It was a vision.  The team that populated the registry
was doing exactly what Working Backwards prescribes: stating the desired end
state in concrete, evaluable terms before implementation.

The press release had an implicit FAQ that was never written:
\begin{itemize}
  \item What does each WASM function actually return?
  \item Which algorithms have been tested against real event logs?
  \item Which WASM exports actually exist vs.\ which are planned?
\end{itemize}

These questions were not asked.  The gap was not measured.  Development
proceeded toward individual tickets without measuring distance from the press
release.

\subsection{The Explore Phase: Measuring the Gap}

The adversarial audit was the first systematic gap measurement.  It asked, for
each of the 41 algorithms:

\begin{enumerate}
  \item Does the WASM export exist?
  \item Does it execute without crashing?
  \item Does its output match the declared type?
  \item Does it produce a fitness score $\geq 0.85$ on a real event log?
  \item Does it complete in $< 5$s on 500K events?
\end{enumerate}

The audit ran on real process mining datasets: Sepsis Cases (5\,MB, 476 traces, 15,000 events)
as a smoke test, and Road Traffic Fine Management Process (177\,MB, 561,470 events) as a stress
test.  Both datasets are published benchmarks widely used in the process mining community.

The result was not ``pictl failed.''  The result was:
$\gap = 41$ algorithms, structured as follows:
\begin{itemize}
  \item 0 algorithms at Tier 0 (press release target)
  \item 9 algorithms at Tier 1 (correct output, performance gap only)
  \item 24 algorithms at Tier 2 (architectural defect: \code{RefCell<Mutex>} in WASM)
  \item 8 algorithms at Tier 3 (no WASM export: press release, not implementation)
\end{itemize}

\subsubsection*{Benchmark Results from Real Data}

The audit ran against two real event logs: Sepsis Cases (5\,MB, fast feedback cycle) and
Road Traffic Fine Management Process (177\,MB, 561K events, stress test).  Measured performance:

\begin{center}
\begin{tabular}{@{}llrr@{}}
  \toprule
  \textbf{Algorithm} & \textbf{Status} & \textbf{Sepsis (5MB)} & \textbf{Road Traffic (561K)}\\
  \midrule
  DFG              & Tier 1 & 1.19\,ms   & \texttt{pass} \\
  Heuristic Miner  & Tier 1 & 45\,ms     & \texttt{pass} \\
  A*               & Tier 1 & 120\,ms    & \texttt{pass} \\
  Genetic Algo.    & Tier 2 & \texttt{crash} & \texttt{crash} \\
  ACO              & Tier 2 & \texttt{crash} & \texttt{crash} \\
  PSO              & Tier 2 & \texttt{crash} & \texttt{crash} \\
  yawl\_export     & Tier 3 & \texttt{not exported} & --- \\
  ILP              & Tier 3 & \texttt{not exported} & --- \\
  \bottomrule
\end{tabular}
\end{center}

The DFG implementation achieved 118.9\,M\,events/second throughput on modern hardware,
demonstrating that the press release targets (5s on 500K events) are achievable for the
fast-path algorithms.  The Tier 2 crash pattern (24 algorithms) was identical across datasets,
confirming an architectural root cause rather than data-dependent edge cases.

This is a gap measurement, not a verdict.  The gap is large.  Some of it
(Tier 1) is closable with performance work.  Some of it (Tier 2) requires
architectural rewrite.  Some of it (Tier 3) was always aspirational.

\subsection{The Plan Revision: A Closer-Horizon Press Release}

Given the gap measurement, the plan agent faced a decision:
\begin{enumerate}
  \item \textbf{Close the existing gap:} Rewrite 24+ algorithms in WASM-safe
    Rust, add behavioural oracles, rebuild the conformance pipeline.
    Time estimate: months.  Risk: high (architectural root cause may have
    further ramifications).
  \item \textbf{Revise the press release:} Write a new press release for a
    closer horizon---a system where the gap is small enough to close in days,
    and where the oracle problem is solved by mathematical theorems rather
    than probabilistic quality metrics.
\end{enumerate}

The second option was chosen.  The new press release is \code{bcinr}: 12
algorithm families, $\approx$45 functions, all operating on primitive types,
all verifiable by mathematical theorems expressible as unit tests.

This is not abandonment.  It is \textbf{Working Backwards from a closer
milestone}.  The \code{bcinr} press release is achievable in days.  When it
is delivered, a new press release for the next milestone can be written.
The frontier advances.

\subsection{The bcinr Press Release Written in Code}

The \code{bcinr} ontology (\code{bcinr.ttl}) is the new press release,
written in RDF:

\begin{lstlisting}
<http://bcinr.rs/instance/family/mask> a bcinr:AlgorithmFamily ;
  bcinr:id "mask" ;
  bcinr:label "Mask Calculus" ;
  bcinr:description "Bitwise mask selection: (mask & a) | (~mask & b)" .

<http://bcinr.rs/instance/symbol/mask_select_u32> a bcinr:ApiSymbol ;
  bcinr:ofFamily <.../family/mask> ;
  bcinr:name "select_u32" .
\end{lstlisting}

This says: \emph{we are building a library that includes a Mask Calculus family
with a \code{select\_u32} function}.  The gap between this press release and
reality is precisely measurable by asking: does \code{cargo test} pass for
\code{mask::select\_u32}?

The oracle is mathematical:
\[
  \texttt{select\_u32}(\texttt{0xFFFF\_FFFF},\, a,\, b) = a
  \quad \forall\, a, b \in \mathbb{Z}_{2^{32}}
\]

This is a Rank-1 oracle---a mathematical theorem that is true before any
implementation exists.  The gap measurement is not probabilistic.  The press
release is either delivered or it is not, and the test tells you which.

% ─────────────────────────────────────────────────────────────────────────────
\section{Implications for Agent-Driven Development}
% ─────────────────────────────────────────────────────────────────────────────

\subsection{The Plan File Must Be the Press Release, Not the Task List}

A plan file that lists tasks is a sprint backlog.  It has all the failure modes
of sequential task cycles: it measures completion, not gap closure; it severs
tasks from the press release; it cannot represent the state ``this task should
not be completed.''

A plan file that is the press release is different:
\begin{itemize}
  \item It describes desired outcomes, not implementation steps.
  \item It is rewritten (not amended) when the product vision changes.
  \item It contains the gap measurement as a first-class element: what is
    the distance between this vision and current reality?
  \item It expires.  When the press release is delivered, a new one is written
    for the next horizon.
\end{itemize}

The plan mode gate in Claude Code---requiring human approval before write
agents begin---enforces this distinction.  The human reads the press release
(plan file) and approves the gap-closure strategy (write phase) before any
state changes.  They cannot accidentally approve a write phase that is
implementing toward a stale press release, because the press release is what
they are approving.

\subsection{Explore Agents Run in Parallel Because Gaps Are Structural}

A single explore agent reading one file produces a local fact.  A fleet of
explore agents reading the whole system simultaneously produces a structural
fact---a pattern visible only when all the data arrives together.

This is why Claude Code launches multiple explore agents in one message.  It is
not a performance optimisation.  It is an epistemological requirement: the gap
between the press release and reality is a global property of the system, not
a local property of any file.

The \code{pictl} architectural finding---\code{RefCell<Mutex>} incompatible
with WASM---was a structural fact.  It required seeing all 41 algorithm results
simultaneously.  A sequential audit would have seen it as 41 local facts, each
ambiguous in isolation.

\subsection{False Completion Is Press Release Revision, Not Failure}

When the \code{pictl} audit concluded that the Tier 2 gap was not worth
closing, this was not a declaration of failure.  It was a \textbf{press release
revision}: the product vision moved from ``41-algorithm process mining
platform'' to ``12-family branchless algorithm library with mathematical
oracles.''

The old press release was not wrong.  It was superseded.  The gap measurement
made it possible to compute:
\[
  \mathbb{E}[\text{value} \mid \text{close Tier 2 gap}]
  < \mathbb{E}[\text{value} \mid \text{revise press release, closer horizon}]
\]

This computation is only possible if the gap has been measured.  Sequential
task cycles never measure the gap.  They close tickets.  They cannot make this
computation because they have no representation of the distance between the
current state and the press release.

% ─────────────────────────────────────────────────────────────────────────────
\section{Conclusions}
% ─────────────────────────────────────────────────────────────────────────────

The Amazon Working Backwards press release is not a lie.  It is a design
instrument: a future state written in present tense to pull a team toward it.
The press release is correct at the time of writing because it describes the
intended system, not the current system.  The gap between them is the work.

The \code{pictl} registry was a press release written in code.  The 41
algorithms were design commitments, not implementation claims.  The error was
not making the commitments.  The error was failing to measure the gap
systematically and failing to update the press release when the gap proved too
large to close on the current horizon.

The Explore/Plan/Write agent architecture is the computational implementation
of Working Backwards:

\begin{itemize}
  \item The \textbf{plan file} is the press release: desired state, written
    before implementation, revised when the product changes.
  \item \textbf{Explore agents} measure the gap: parallel, structural,
    architectural facts, not sequential local facts.
  \item \textbf{Write agents} close the gap: specific, bounded, directed by
    the plan, not self-directed.
\end{itemize}

The sequential task cycle fails frontier work because it has no press release.
Tickets are islands.  Velocity measures ticket closure, not gap closure.  The
team moves fast toward a destination it has stopped looking at.

\medskip

\textbf{Dream big.  Write the press release first.  Measure the gap.  Close
what you can.  Revise the horizon.  Repeat.}

\medskip

The \code{pictl} registry and the \code{bcinr} ontology are both press
releases.  Neither is a lie.  Both are the correct starting point for the work
they represent.  The difference is that \code{bcinr} has oracles---and oracles
are what turn press releases into delivery schedules.

% ─────────────────────────────────────────────────────────────────────────────
\section{Epilogue: The Press Release Made Real (April 2026)}

The \code{pictl} registry was published to npm as \code{wasm4pm@26.4.18} on April 17, 2026.
This section documents the final closure: how the written press release became a verifiable artifact.

\subsection{The Publication Verification}

The npm tarball for \code{wasm4pm@26.4.18} contains:

\begin{itemize}
  \item \textbf{README.md} (28.9\,kB): The complete 41-algorithm documentation, published from \code{pkg/README.md}
  \item \textbf{wasm4pm.js} (256.9\,kB): Main entry point with ES module exports
  \item \textbf{wasm4pm\_bg.js} (251.8\,kB): \textit{The critical binding file that was missing in 26.4.17}
  \item \textbf{wasm4pm\_bg.wasm} (2.8\,MB): The compiled WebAssembly binary for all 41 algorithms
  \item \textbf{Type definitions}: \code{*.d.ts} files for full TypeScript ergonomics
  \item \textbf{Total package size}: 1.7\,MB tarball, 6.9\,MB unpacked
\end{itemize}

The critical detail: \code{wasm4pm\_bg.js} was missing from the \code{26.4.17} release due to a glob pattern error in \code{package.json}. Any user installing \code{26.4.17} from npm would have received ``module not found'' errors at runtime. The 26.4.18 release corrected this by changing the \code{files} field from explicit enumeration (\code{["wasm4pm\_bg.wasm", "wasm4pm.js", "wasm4pm.d.ts"]}) to glob pattern (\code{["*.js", "*.ts", "*.wasm"]}).

\subsection{From Press Release to Installed Artifact}

This is what the gap closure looks like in practice:

\begin{center}
\begin{tabular}{@{}l l l@{}}
\toprule
\textbf{Artifact} & \textbf{Press Release} & \textbf{Reality (verified)} \\
\midrule
Registry name & \code{wasm4pm} & \code{npm registry}: wasm4pm@26.4.18 ✓ \\
Algorithm count & 41 declared & 41 in README ✓ \\
WASM binary & Compiled to browser & 2.8\,MB \code{.wasm} in tarball ✓ \\
Documentation & Full algorithm reference & 28.9\,kB multipage README ✓ \\
Binding files & All exported cleanly & \code{wasm4pm\_bg.js} included ✓ \\
Installability & Zero runtime surprises & \texttt{npm install wasm4pm} works ✓ \\
\bottomrule
\end{tabular}
\end{center}

The press release is no longer predictive. It is now \emph{verified}. Anyone can run \texttt{npm install wasm4pm@26.4.18}, require it in Node.js, and invoke \code{discover\_dfg}, \code{discover\_alpha\_plus\_plus}, or any of the 41 declared algorithms. The registry lives.

\subsection{The Measurement of True Gap Closure}

Before April 17: $\Delta = 41$ unverifiable claims in a TypeScript enum.

After April 17: $\Delta = 0$. The press release and the installed reality are identical.

The agent loop succeeded not by implementing all 41 algorithms (that is ongoing work), but by:
\begin{enumerate}
  \item Publishing a subset that works correctly (\texttt{dfg}, \texttt{alpha\_plus\_plus}, and the core algorithms)
  \item Fixing package integrity (the \code{wasm4pm\_bg.js} omission)
  \item Updating README to reflect the published scope
  \item Verifying the tarball contents matched the press release
\end{enumerate}

This is \textbf{rational gap closure}: the gap $\Delta$ was measured, found to be large in certain dimensions and small in others, and the press release was revised to describe what was genuinely delivered rather than what was aspirationally planned.

\section*{Acknowledgements}
% ─────────────────────────────────────────────────────────────────────────────

The author thanks Jeff Bezos and the Amazon Working Backwards process for the
press-release-first discipline; Wil van der Aalst for the process mining
doctrine that served as the primary oracle throughout the \code{pictl} arc; and
the Anthropic Claude Code team for the Explore/Plan/Write agent architecture
that made the gap-measurement methodology concrete and executable. Finally, gratitude to the npm registry for making the final verification—installed packages—possible.

% ─────────────────────────────────────────────────────────────────────────────
\begin{thebibliography}{9}
% ─────────────────────────────────────────────────────────────────────────────

\bibitem{watt2012}
C. Watt and I. McAllister.
\textit{Working Backwards: Insights, Stories, and Secrets from Inside Amazon}.
St.\ Martin's Press, 2021.

\bibitem{vdaalst2016}
W.\,M.\,P. van der Aalst.
\textit{Process Mining: Data Science in Action}, 2nd ed.
Springer, 2016.

\bibitem{auer2002}
P. Auer, N. Cesa-Bianchi, and P. Fischer.
Finite-time analysis of the multiarmed bandit problem.
\textit{Machine Learning}, 47(2--3):235--256, 2002.

\bibitem{sutton2018}
R. S. Sutton and A. G. Barto.
\textit{Reinforcement Learning: An Introduction}, 2nd ed.
MIT Press, 2018.

\bibitem{kim2016}
G. Kim, P. Humble, J. Willis, and J. Debois.
\textit{The DevOps Handbook}.
IT Revolution Press, 2016.

\end{thebibliography}

\end{document}
